\documentclass[12pt]{article}
\usepackage[english]{babel}
\usepackage[style=authoryear,backend=biber,maxcitenames=2]{biblatex}
\usepackage{csquotes}
\usepackage{setspace}
\usepackage[a4paper,margin=2.5cm]{geometry}
\usepackage{footmisc}
\addbibresource{refs.bib}
\doublespacing

\title{On Reading Slowly}
\author{Eleanor Vance}
\date{}

\begin{document}
\maketitle

\section*{Introduction}
As \textcite{smith2020} argue, \enquote{reading is a practice, not a product}.
Others disagree \parencite[p.~12]{knuth1984}, and some go further \autocite{lee2019}.
\citeauthor{smith2020} published in \citeyear{smith2020}. A footnote citation is common.\footcite{lee2019}
Multiple works: \parencite{smith2020,knuth1984}; see also \cite[see][ch. 2]{knuth1984}.
\Textcite{lee2019} is often cited. \citetitle{knuth1984} remains influential.

\begin{displayquote}
Reading slowly is an act of resistance against the speed of the attention economy.
\end{displayquote}

The phrase \enquote*{close reading} has a long history.\footnote{See the discussion in \cite{smith2020}.}

\begin{quote}
\textit{``Festina lente''} --- make haste slowly.
\end{quote}

\nocite{web2023}
\printbibliography[title={Works Cited}, heading=bibintoc]
\end{document}
